\section{Related Work}
\label{sec:related}

A line of recent work asks whether a video generation model can itself be the reasoner: given a first frame and a prompt, the model renders the solution as a video rather than describing it. \citet{wiedemer2025videomodels} document such zero-shot ``chain-of-frames'' behaviour in Veo~3 across perception, manipulation and maze tasks, and \citet{tong2025thinkingwithvideo} frame video generation as a general multimodal reasoning paradigm. VBVR \citep{vbvrpro2026} and its successor VBVR-Pro \citep{vbvrpro2026} make the question measurable at scale. Hundreds of procedural generators produce first-frame, prompt and answer-video triples, a rule-based evaluator scores each family, and VBVR-Pro-Bench holds out 50 families from all training to separate memorisation from generalisation. VBVR-Pro also post-trains video models with the same evaluators as verifiable rewards, following GRPO \citep{shao2024deepseekmath,guo2025deepseekr1} as adapted to flow and diffusion models by Flow-GRPO \citep{liu2025flowgrpo} and DanceGRPO \citep{xue2025dancegrpo}. Its RL-trained checkpoints are the strongest video models on the benchmark and the natural ceiling for the comparison in this paper. We keep the benchmark, the evaluator and the published video-model numbers untouched and change only the family of system under test.

Writing a program is an older route to visual reasoning. Visual Programming \citep{gupta2023visprog} and ViperGPT \citep{suris2023vipergpt} have a language model compose vision modules into an executable program, and Visual Sketchpad \citep{hu2024sketchpad} lets a multimodal model draw intermediate results with code. These systems answer questions in text. Ours must render a full video, which turns the program from a scaffold into the deliverable. On the language side, program-aided reasoning \citep{gao2023pal,chen2023pot} and Code as Policies \citep{liang2023codeaspolicies} showed that offloading computation to an interpreter improves reliability. The agents we evaluate are the current form of that idea, Codex-style models \citep{chen2021codex} wrapped in a terminal loop that reads files, runs commands and iterates on errors. SWE-bench and SWE-agent \citep{jimenez2024swebench,yang2024sweagent} established the evaluation pattern we borrow (one sandboxed container per instance, one attempt, a verifier decides), and the terminal-agent benchmark tradition \citep{harbor2025} packages tasks the same way. We are not aware of prior work that runs off-the-shelf coding agents, unmodified, against a video-reasoning leaderboard.

Video benchmarks mostly score with learned components. VBench \citep{huang2024vbench} mixes feature-based metrics, VideoScore \citep{he2024videoscore} trains a video-language judge, and LLM-as-a-judge protocols inherit known position and verbosity biases \citep{zheng2023judging}. Physics-IQ \citep{motamed2025physicsiq} and VBVR-Pro instead compute scores from the pixels with classical computer vision against symbolic ground truth, so the same output always receives the same score and the score can serve as a reward. The same property makes the comparison in this paper exact, since a coding agent and a video model are graded by one deterministic function on the same 500 instances.
